% !PRINTED 512
\weberpage{512}

1. If the conjugate groups $\pi^{-1}Q\pi$ are coprime, then
the simultaneous adjunction of all the roots of the
resolvent $\varphi(t) = 0$ is (by Theorem 4, \S.\,155) equivalent to
the adjunction of a function belonging to the unit group,
and the solution of the given equation is reduced to the complete
solution of the resolvent. We have
\[ N = \Omega(\psi, \psi_1 \ldots \psi_{j-1}), \]
and a \gesperrt{Galois} resolvent of the equation $\varphi(t) = 0$ is
at the same time a \gesperrt{Galois} resolvent of the original equation.
In this case we call $\varphi(t) = 0$ a \gesperrt{total resolvent}
of the given equation.

2. If the conjugate groups $\pi^{-1}Q\pi$ have a divisor $R$
of degree $r$ different from the unit group, which is then
a normal divisor of $P$, then the simultaneous adjunction
of all the functions conjugate to $\psi$ is equivalent
to the adjunction of a quantity belonging to $R$. The \gesperrt{Galois}
resolvent of the given equation is not yet completely solved
by this adjunction; rather it is only decomposed into factors
of degree $r$. In this case $\varphi(t) = 0$ is called a
\gesperrt{partial resolvent}.

If $P$ is a \gesperrt{simple group}, then \gesperrt{only} total resolvents
exist, whereas, if $P$ has normal divisors, a partial
resolvent can be found for every such normal
divisor.

The same distinction also comes to light when we investigate the
\gesperrt{Galois} group of the resolvent $\varphi(t)$. This
group consists of all the permutations that are produced in the sequence of the
quantities
\weq{1}{\psi, \psi_1, \psi_2 \ldots \psi_{j-1}}
by the application of all the operations $\pi$ of $P$; for every equation in $\Omega$ between the quantities (1)
remains valid when such a permutation is carried out,
since one can apply the operations $\pi$ to such an equation;
and if a function in $\Omega$ of the quantities (1) admits all these
permutations, then it also admits all the permutations $\pi$
and is therefore equal to a quantity in $\Omega$.

The degree of this group is now to be determined. Among
the operations of $P$, those and only those will produce no
change among the quantities (1) which occur simultaneously in
%CONT>>
